\documentclass[11pt, a4paper]{article}

% ── Packages ──────────────────────────────────────────────────────────────────
\usepackage[top=2.2cm, bottom=2.2cm, left=2.5cm, right=2.5cm]{geometry}
\usepackage{graphicx}
\usepackage{booktabs}
\usepackage{array}
\usepackage{caption}
\usepackage{subcaption}
\usepackage{hyperref}
\usepackage{parskip}
% \usepackage{microtype} % disabled
\usepackage{titlesec}
% \usepackage{lmodern} % replaced by mlmodern
\usepackage[T1]{fontenc}
\usepackage[utf8]{inputenc}
\usepackage{xcolor}
\usepackage{enumitem}
\usepackage{amsmath}
\usepackage{fancyhdr}
\usepackage{lastpage}

% ── Hyperref setup ────────────────────────────────────────────────────────────
\hypersetup{
    colorlinks = true,
    linkcolor  = blue!60!black,
    urlcolor   = blue!60!black,
    citecolor  = blue!60!black,
    pdftitle   = {Tradable Prices and FDI – Argentina and Costa Rica},
    pdfauthor  = {Camila Aguirre},
}

% ── Section formatting ────────────────────────────────────────────────────────
\titleformat{\section}
  {\large\bfseries\color{black!80}}{\thesection.}{0.5em}{}[\titlerule]
\titleformat{\subsection}
  {\normalsize\bfseries}{\thesubsection.}{0.5em}{}

\titlespacing{\section}{0pt}{14pt}{6pt}
\titlespacing{\subsection}{0pt}{10pt}{4pt}

% ── Header / Footer ───────────────────────────────────────────────────────────
\pagestyle{fancy}
\fancyhf{}
\renewcommand{\headrulewidth}{0.4pt}
\fancyhead[L]{\small\textit{Tradable Prices \& FDI — Argentina and Costa Rica}}
\fancyhead[R]{\small\textit{\today}}
\fancyfoot[C]{\small Page \thepage\ of \pageref{LastPage}}

% ── Caption setup ─────────────────────────────────────────────────────────────
\captionsetup{
    font       = small,
    labelfont  = bf,
    margin     = 10pt,
    skip       = 6pt,
}

% ── Misc ──────────────────────────────────────────────────────────────────────
\setlength{\parindent}{0pt}
\setlength{\parskip}{6pt}
\renewcommand{\arraystretch}{1.3}

% ══════════════════════════════════════════════════════════════════════════════
\begin{document}
% ══════════════════════════════════════════════════════════════════════════════

% ── Title block ───────────────────────────────────────────────────────────────
\begin{center}
    {\LARGE\bfseries Tradable and Non-Tradable Consumer Prices\\[4pt]
     and Foreign Direct Investment}\\[8pt]
    {\large Argentina (Part 1) \quad|\quad Costa Rica (Part 2)}\\[6pt]
    {\large\bfseries Camila Aguirre}\\[4pt]
    {\normalsize\color{gray}\today}
\end{center}

\vspace{4pt}
\hrule
\vspace{10pt}

% ══════════════════════════════════════════════════════════════════════════════
\section*{Nomenclature}
% ══════════════════════════════════════════════════════════════════════════════

\begin{tabular}{>{\bfseries}p{3.4cm} p{10.8cm}}
    \toprule
    \normalfont\textbf{Acronym} & \textbf{Definition} \\
    \midrule
    \multicolumn{2}{l}{\textit{Part 1 — Argentina}} \\[3pt]
    CPI          & Consumer Price Index \\
    IPC          & \emph{Índice de Precios al Consumidor} \\
    IPC-GBA      & IPC Gran Buenos Aires \\
    IPC Nacional & National CPI, Total Nacional \\
    Bienes       & Goods component of the IPC (tradable proxy) \\
    Servicios    & Services component of the IPC (non-tradable proxy) \\
    YoY          & Year-on-year \\
    INDEC        & \emph{Instituto Nacional de Estadística y Censos} \\[6pt]
    \multicolumn{2}{l}{\textit{Part 2 — Costa Rica}} \\[3pt]
    FDI          & Foreign Direct Investment \\
    IIP          & International Investment Position \\
    BOP          & Balance of Payments \\
    OVAP         & \emph{Otras Variaciones de Activos y Pasivos}
                   (Other Changes in Position) \\
    BPM6         & IMF Balance of Payments and IIP Manual, 6th Edition \\
    BCCR         & \emph{Banco Central de Costa Rica} \\
    IMF          & International Monetary Fund \\
    \bottomrule
\end{tabular}

\vspace{12pt}

% ══════════════════════════════════════════════════════════════════════════════
\section{Part 1 — Argentina: Tradable and Non-Tradable Consumer Prices}
% ══════════════════════════════════════════════════════════════════════════════

\subsection{Data Sources}

All data are sourced from \textbf{INDEC} (Instituto Nacional de Estadística y
Censos), the official Argentine statistical authority
(\url{https://www.indec.gob.ar}).

Argentina's CPI statistics are produced by two surveys that are spliced
to cover the full 2010--present period. The current \textbf{IPC Nacional}
(Total Nacional, base December~2016\,=\,100) was launched in December~2016;
prior to that, the \textbf{IPC-GBA} (Greater Buenos Aires, same base) is
used as the bridging source for April–November~2016. Before April~2016 no
official series with this methodology is available, so a statistical
backcast is applied (see below).

The two files downloaded at runtime are:

\begin{itemize}[leftmargin=1.5em, itemsep=2pt]
    \item \texttt{sh\_ipc\_aperturas.xls}, sheet \emph{Índices aperturas}:
          IPC Nacional, base December~2016\,=\,100, December~2016 to present.

    \item \texttt{sh\_ipc\_12\_16.xls}: IPC-GBA, base December~2016\,=\,100,
          April–November~2016. Used as the bridge to the backcast period.
\end{itemize}

Both files are downloaded programmatically at runtime via \texttt{requests};
no manual file handling is required.

\subsection{Classification: Tradable vs.\ Non-Tradable}

The Bienes / Servicios split published by INDEC is used directly as the
tradable / non-tradable classification, following the
\textbf{Balassa--Samuelson framework}:

\begin{itemize}[leftmargin=1.5em, itemsep=2pt]
    \item \textbf{Tradable} $\equiv$ \emph{Bienes} (goods): prices of goods
          are arbitraged across borders through international trade and
          exchange-rate movements. The categories included are Food \&
          non-alcoholic beverages (20.40\%), Alcoholic beverages \& tobacco
          (3.09\%), Clothing \& footwear (5.53\%), and Household equipment
          (5.49\%), with a combined INDEC weight of 34.51\%.

    \item \textbf{Non-Tradable} $\equiv$ \emph{Servicios} (services): service
          prices are predominantly determined by domestic costs (wages, rents)
          and face limited international competition. The categories included
          are Housing \& utilities (14.01\%), Health (9.51\%), Transport
          (11.34\%), Communication (3.36\%), Recreation (7.68\%), Education
          (6.46\%), Restaurants \& hotels (8.47\%), and Miscellaneous
          (4.66\%), with a combined weight of 65.49\%.
\end{itemize}

\textit{Note}: The residual categories not assigned to either group
(primarily those with mixed tradability) represent a minor share of the
basket and are excluded from the incidence calculation; overall CPI
(\emph{Nivel General}) is retained as the reference total.

\subsection{Data Processing and Construction}

\textbf{Step 1 — Dec~2016–present (IPC Nacional).}
The \emph{Bienes} and \emph{Servicios} index rows are read directly from
\texttt{sh\_ipc\_aperturas.xls}. If the direct rows are absent, a weighted
average over the constituent divisions is computed using the INDEC weights
listed above.

\textbf{Step 2 — Apr–Nov~2016 (IPC-GBA bridge).}
The \emph{Bienes} and \emph{Servicios} rows are read from
\texttt{sh\_ipc\_12\_16.xls}. Both surveys share December~2016\,=\,100 as
their base, so the two segments are \textbf{spliced directly at the level
without any rescaling}.

\textbf{Step 3 — Backcast Jan~2010–Mar~2016.}

\begin{center}
\colorbox{yellow!20}{%
\begin{minipage}{0.93\textwidth}
\small
\textbf{\color{red!70!black}Data Limitation — Pre-2016 period.}
No official INDEC series with the December~2016 base extends before
April~2016. The Jan~2010--Mar~2016 segment is a \textbf{statistical
backcast}, not observed data, and should be \textbf{used with caution}.
It is shaded in all charts to signal this limitation.
\end{minipage}}
\end{center}

\vspace{4pt}

A \textbf{log-linear trend} is fitted by OLS on the April–November~2016
window (8~observations):
\[
    \ln P_t \;=\; a + b\,t \;+\; \varepsilon_t
\]
The estimated monthly growth rate $\hat{b}$ implies annual rates of
approximately 25–35\% for \emph{Bienes} and 20–30\% for \emph{Servicios},
consistent with third-party inflation estimates for Argentina in 2010--2016.
The backcast is:
\[
    \hat{P}_{t-k} \;=\; P_{\text{Apr\,2016}} \cdot e^{-\hat{b}\,k}
\]

\textbf{Step 4 — Rebase to January~2010\,=\,100.}
The three segments are concatenated, de-duplicated, and rescaled so that
January~2010\,=\,100 for all series.

\textbf{Step 5 — Derived indicators.}
Year-on-year (YoY) inflation is the 12-month percentage change of the index.
Monthly incidence is:
\[
    \text{Inc}_{i,t} \;=\; w_i \times \Delta\%P_{i,t}
\]
where $w_i$ is the INDEC base-Dec-2016 weight and $\Delta\%P_{i,t}$ the
monthly percentage change of component~$i$.

\subsection{Results}

\begin{figure}[ht]
    \centering
    \includegraphics[width=\textwidth]{fig1_yoy_inflation.pdf}
    \caption{%
        \textbf{Year-on-year CPI inflation — Tradable vs.\ Non-Tradable
        (Argentina, 2010–present).}
        The orange line shows annual inflation for \emph{Bienes} (tradable
        goods); the blue dashed line shows \emph{Servicios} (non-tradable
        services). Both series are indexed to January~2010\,=\,100 and
        YoY changes are computed as 12-month percentage changes.
        The shaded area (2010–2016) corresponds to the log-linear
        backcast; the vertical dotted line marks the Dec~2016 splice
        between IPC-GBA and IPC Nacional.
        \textit{Source: INDEC (\texttt{sh\_ipc\_aperturas.xls},
        \texttt{sh\_ipc\_12\_16.xls}).}
    }
    \label{fig:yoy}
\end{figure}

\begin{figure}[ht]
    \centering
    \includegraphics[width=\textwidth]{fig2_monthly_incidence.pdf}
    \caption{%
        \textbf{Monthly CPI incidence by component — Tradable vs.\
        Non-Tradable (Argentina, recent 5~years).}
        Stacked bars decompose the monthly change of the overall CPI
        (\emph{Nivel General}) into the contribution of \emph{Bienes}
        (tradable, orange) and \emph{Servicios} (non-tradable, blue).
        The black line shows the total monthly CPI change for reference.
        Incidence is computed as $w_i \times \Delta\%P_{i,t}$ using
        fixed INDEC weights (base December~2016).
        Positive bars indicate upward price pressure; negative bars
        indicate downward contributions in that month.
        \textit{Source: INDEC (\texttt{sh\_ipc\_aperturas.xls}).}
    }
    \label{fig:incidence}
\end{figure}

% ══════════════════════════════════════════════════════════════════════════════
\section{Part 2 — Costa Rica: FDI, IIP vs.\ Balance of Payments}
% ══════════════════════════════════════════════════════════════════════════════

\subsection{Data Sources}

\textbf{Primary source}: IMF Balance of Payments and International
Investment Position Statistics, accessed via the public
\textbf{IMF SDMX-JSON REST API}
(\url{https://data.imf.org/en/Resource-Pages/IMF-API}).
No authentication is required. The dataset covers over 190~economies
under the BPM6 framework.

\textbf{Fallback source}: If the API is unreachable at runtime, the
script uses hardcoded values transcribed from the
\textbf{Banco Central de Costa Rica (BCCR)} official statistical portal,
accessed on 25~May~2025. The BCCR is
the national compiler of BOP and IIP statistics for Costa Rica and
submits these figures to the IMF; the two sources are therefore
consistent.

The two series collected are:

\begin{center}
\begin{tabular}{>{\ttfamily}l l l}
    \toprule
    \normalfont\textbf{Indicator} & \textbf{Description} & \textbf{Sign convention} \\
    \midrule
    BFDI\_BP6\_USD & Financial acct., direct investment,
                     net incurrence of liabilities & $+$ = net inflow \\
    ILFD\_BP6\_USD & IIP liabilities, direct investment
                     (end-of-period stock)         & $+$ = outstanding stock \\
    \bottomrule
\end{tabular}
\end{center}

Both indicators measure \emph{inward} FDI from the liability side of
Costa Rica's balance sheet and share the same BPM6 sign convention.

\subsection{Aggregation from Quarterly to Annual}

The IMF BOP dataset publishes data at \textbf{quarterly} frequency.
Annual series are constructed as follows:

\begin{itemize}[leftmargin=1.5em, itemsep=2pt]
    \item \textbf{BOP flows} (\texttt{BFDI\_BP6\_USD}):
          Annual\,=\,\textbf{sum} of Q1\,+\,Q2\,+\,Q3\,+\,Q4.
          \textit{Justification}: flows are additive across time
          sub-periods — a transaction recorded in Q2 is economically
          equivalent to one in Q4 (BPM6~§3.7). Summing gives total
          transactions during the calendar year.

    \item \textbf{IIP stock} (\texttt{ILFD\_BP6\_USD}):
          Annual\,=\,\textbf{Q4 end-of-year value} (31~December).
          \textit{Justification}: stocks represent a balance at a
          specific point in time, not cumulative transactions.
          Averaging or summing quarterly stocks would be conceptually
          incorrect (BPM6~§7.4). The year-end (Q4) value is the
          standard annual IIP benchmark.
\end{itemize}

\subsection{Sign Convention Verification}

The script verifies empirically that both series use compatible sign
conventions by checking three conditions: (i)~the majority of annual
BOP flow observations are positive (net inflows); (ii)~the IIP stock
is consistently positive and generally increasing; and
(iii)~the correlation between cumulated BOP flows and the IIP stock
exceeds 0.90. For Costa Rica, all three conditions are satisfied
(correlation\,$\approx$\,0.997), confirming that \textbf{no sign
adjustment is required}.

\subsection{Indicator Construction}

The BPM6 reconciliation identity (Chapter~9) states:
\[
    \text{IIP}_t \;=\; \text{IIP}_{t-1}
                 \;+\; \text{Flow}_t
                 \;+\; \text{OtherChanges}_t
\]
where \emph{OtherChanges} (OVAP in Spanish) captures valuation
effects (price and exchange-rate changes), reclassifications, and
statistical revisions not recorded as transactions.

The \textbf{annual residual} is:
\[
    \text{Residual}_t \;=\; \Delta\text{IIP}_t \;-\; \text{Flow}_t
\]

\textbf{Index construction.}
To make the IIP stock (a level) and BOP flows (period transactions)
directly comparable on a single chart, both are expressed as index
numbers with base year 2015\,=\,100:
\[
    \text{IIP index}_t
        \;=\; \frac{\text{IIP}_t}{\text{IIP}_{2015}} \times 100
\]
\[
    \text{BOP index}_t
        \;=\; \frac{\text{IIP}_{2015} + \sum_{s=2016}^{t}\text{Flow}_s}
                   {\text{IIP}_{2015}} \times 100
\]
The BOP index is the \emph{flow-implied position}: the level the IIP
stock would reach if only BOP-recorded transactions had driven it,
anchored to the actual 2015 level. The vertical gap between the two
lines equals the accumulated OVAP residual.

\clearpage
\subsection{Results}

\begin{figure}[ht]
    \centering
    \includegraphics[width=\textwidth]{figure1_fdi_index.pdf}
    \caption{%
        \textbf{Inward FDI — IIP position vs.\ cumulated BOP flows,
        index 2015\,=\,100 (Costa Rica, 2015–2024).}
        The blue line (\emph{IIP index}) shows how the end-of-year
        inward FDI stock has evolved relative to 2015.
        The red dashed line (\emph{BOP cumulated-flows index}) shows
        what the stock would be if only BOP-recorded transactions had
        driven it (anchored to the 2015 IIP level).
        The shaded gap between the two lines represents the
        accumulated Other Changes in Position (OVAP): valuation effects,
        exchange-rate changes, reclassifications, and statistical
        revisions. A positive gap indicates the IIP stock grew faster
        than transactions alone would imply.
        \textit{Source: IMF BOP/IIP dataset — SDMX-JSON API
        (\texttt{BFDI\_BP6\_USD}, \texttt{ILFD\_BP6\_USD}); fallback:
        BCCR. Framework: BPM6.}
    }
    \label{fig:fdi_index}
\end{figure}

\begin{figure}[ht]
    \centering
    \includegraphics[width=\textwidth]{figure2_fdi_residual.pdf}
    \caption{%
        \textbf{Other Changes in Position (OVAP / Residual):
        $\Delta\text{IIP}_t - \text{Flow}_t$
        (Costa Rica, 2016–2024).}
        Each bar shows the annual residual between the change in the
        IIP position and the net FDI flow recorded in the Balance of
        Payments. A \textbf{blue} (positive) bar indicates that the
        IIP stock grew \emph{more} than the recorded flow — consistent
        with valuation gains or upward statistical revisions.
        A \textbf{red} (negative) bar indicates the reverse.
        The annotation box reports the cumulative residual over the
        full period and its share of total $\Delta\text{IIP}$.
        \textit{Source: IMF BOP/IIP dataset — SDMX-JSON API
        (\texttt{BFDI\_BP6\_USD}, \texttt{ILFD\_BP6\_USD}); fallback:
        BCCR. Framework: BPM6, Chapter~9.}
    }
    \label{fig:fdi_residual}
\end{figure}

% ── Footer note ───────────────────────────────────────────────────────────────
\vfill
\begin{center}
    \small\color{gray}
    All Python scripts used to produce the figures are provided
    alongside this report and are executable end-to-end without
    manual intervention.
\end{center}

\end{document}
